We give an outline of the derivation of necessary conditions for the existence of Turing patterns given in \cite{murray03}. Consider the concentrations of two chemical substances. Their dynamics are modelled by two coupled reaction-diffusion equations. Let $\Omega \subseteq \R^n$ be an open and bounded set with $C^1$ boundary and $T > 0$. Let $u(t, x)$ and $v(t, x)$ denote the concentrations of the two species at $t \in [0, T]$ and $x \in \overline{\Omega}$. A reaction-diffusion model for these concentrations $u$ and $v$ is of the form
%
\begin{equation}
    \label{eq:reac_diff}
    \begin{aligned}
        u_t &= \Delta u + \gamma f(u, v) & \text{in } &(0, T) \times \Omega, \\
        v_t &= d \Delta v + \gamma g(u, v) & \text{in } &(0, T) \times \Omega, \\
        \pderiv[u]{\nu} &= \pderiv[v]{\nu} = 0 & \text{on } &[0, T) \times \partial \Omega, \\
        u(0, \cdot\ ) &= u_0 & \text{in } &\Omega, \\
        v(0, \cdot\ ) &= v_0 & \text{in } &\Omega,
    \end{aligned}
\end{equation}

\noindent
where $f: \R^2 \to \R$ and $g: \R^2 \to \R$ are the reaction functions, $d > 0$ is the diffusion coefficient for $v$, $\gamma > 0$ is a reaction scaling coefficient and $u_0, v_0: \Omega \to \R$ are given initial data. Turing patterns can emerge when this type of system has a homogeneous equilibrium that is linearly stable in the absence of diffusion, i.e. if it is a linearly stable steady state of the system
%
\begin{equation}
    \label{eq:reac_nodiff}
    u_t = \gamma f(u, v), \qquad v_t = \gamma g(u, v),
\end{equation}

\noindent
and becomes unstable due to the introduction of diffusion in \cref{eq:reac_diff}, referred to as a \emph{diffusion driven instability}. The relevant steady state is the positive solution $(U, V)$ of
%
\begin{equation*}
    f(U, V) = g(U, V) = 0,
\end{equation*}

\noindent
if it exists. To first investigate the stability of the steady state in absence of diffusion, we linearise the system of \cref{eq:reac_nodiff} around $(U, V)$ and compute the eigenvalues of the corresponding matrix. For the steady state to be asymptotically stable, we require that $\Re{\lambda} < 0$ for every eigenvalue $\lambda \in \mathbb{C}$ of this matrix, for which we should have
%
\begin{equation}
    \label{eq:pattern_conds_1}
    f_u + g_v < 0 \qquad \text{and} \qquad f_u g_v - f_v g_u > 0,
\end{equation}

\noindent
where all partial derivatives are evaluated at the steady state $(U, V)$. To study the instability of the steady state with respect to the system of \cref{eq:reac_diff}, the system is linearised and decomposed into \emph{Fourier modes} $w_k$ with \emph{wave number} $k$\footnote{These Fourier nodes and wave numbers are the solutions of the problem $\Delta w_k + k^2 w_k = 0$. They depend on the dimension $n$, the shape of the spatial domain $\Omega$ as well as the boundary conditions of \cref{eq:reac_diff}. For more details, see \cite{murray03}.}. We assume that the solution can be written as
%
\begin{equation*}
    \begin{pmatrix}u(t, x) - U \\ v(t, x) - V \end{pmatrix} = w(t, x) = \sum_{k} c_k e^{\lambda(k)t} w_k(x) ,
\end{equation*}

\noindent
for coefficients $c_k \in \R$ and growth rates $\lambda(k) \in \mathbb{C}$ for every wave number $k$. If the growth rate $\lambda(k)$ has positive real part for some wave number $k$, the corresponding Fourier mode grows exponentially in time, causing the steady state to become unstable. Which Fourier modes precisely cause the instability depends on $\gamma$ and influences the spatial scale of the patterns.

One can show that these considerations lead to the following additional set of necessary conditions for the reaction functions $f$ and $g$,
%
\begin{equation}
    \label{eq:pattern_conds_2}
    d f_u + g_v > 0 \qquad \text{and} \qquad \left(d f_u + g_v\right)^2 - 4d \left(f_u g_v - f_v g_u\right) > 0,
\end{equation}

\noindent
where all partial derivatives are evaluated at $(U, V)$. In particular, the conditions of \cref{eq:pattern_conds_1,eq:pattern_conds_2} require $f_u$ and $g_v$ to be of opposite sign, and $d > 1$. Chemically, this is interpreted as one of the two species diffusing faster than the other. Lastly, the reaction scaling coefficient $\gamma$ should be set such that the spatial scale of the patterns is comparable to the size of the spatial domain $\Omega$.